\documentclass[11pt]{article}
\usepackage[a4paper,margin=2.2cm]{geometry}
\usepackage{xcolor}
\usepackage{fontspec}
\usepackage{xeCJK}
\usepackage{enumitem}
\setmainfont{Times New Roman}
\setCJKmainfont{SimSun}
\setlist[itemize]{leftmargin=1.4em,itemsep=0.35em,topsep=0.35em}
\setlength{\parindent}{0pt}
\setlength{\parskip}{0.55em}
\pagestyle{plain}
\begin{document}
{\color[HTML]{1F5FD2}\LARGE\bfseries 辅助沟通系统（AAC）应用指南\par}
{\color[HTML]{667085}\small 星轨原创教育资源：用于家庭与学校共同制定支持计划。\par}
\vspace{0.8em}
{\color[HTML]{111827}\large\bfseries AAC 是什么\par}
\begin{itemize}
\item 辅助沟通系统包括手势、图片、沟通板、书写、语音输出设备和应用程序，目的是让儿童拥有可用的表达通道。
\item AAC 不会阻碍口语发展，反而能在口语尚未稳定时减少挫败，帮助儿童更早参与互动。
\end{itemize}
{\color[HTML]{111827}\large\bfseries 选择系统\par}
\begin{itemize}
\item 选择 AAC 时要综合考虑视觉、动作操作、理解能力、语言背景、家庭流程、学校场景和儿童最需要表达的信息。
\item 系统不宜只包含名词，还应包含请求、拒绝、评论、求助、情绪和社交问候等功能性信息。
\end{itemize}
{\color[HTML]{111827}\large\bfseries 教学策略\par}
\begin{itemize}
\item 成人应先示范 AAC 的使用。例如递饼干时指向“想要”和“饼干”，不要求儿童立即完全正确。
\item 在高动机活动中等待儿童尝试，并把不完整的指点、眼神或声音都当作有意义的沟通信号。
\end{itemize}
\vfill
{\color[HTML]{667085}\footnotesize 本材料为应用内原创教育内容，不能替代医疗、法律或正式评估建议。\par}
\end{document}
